\chapter{Signal efficiency and interference}
\label{app:signal_interference_study}

Signal efficiencies and the effect of signal--SM interference are examined for the regions defined in Chapter~\ref{chapter:event_selection}.

\section{Signal efficiency}
\label{app:signal_efficiency}

Signal efficiencies are evaluated using the pure BSM contribution.
The acceptance times selection efficiency, $A\epsilon$, is evaluated with the total generated signal as its reference and includes the kinematic acceptance and reconstruction and selection losses.
For illustration, at $(\MU,\gU,\beta_L^{23})=(1.5~\TeV,1.0,0.2)$ with $\beta_R^{33}=0$, $A\epsilon$ is $0.30\%$ and $0.13\%$ in $\bVetoSR$-$\Res$ and $\bVetoSR$-$\NonRes$, compared with $0.64\%$ and $0.069\%$ in the corresponding $\bTagSR$ regions.
These efficiencies depend on the mass and flavour couplings because the relative contributions of the production and decay channels change.

\section{Signal--SM interference}

The interference term is defined in Eq.~\eqref{eq:signal_interference}. Its contribution is examined here after the SR selections.
The low-energy phase space in which interference with SM $W$ production is most important is suppressed by the requirements on $\pT^{\hadtau}$, $\met$ and $\mT$.
In particular, the non-resonant requirements $\pT^{\hadtau}>200~\GeV$, $\met>400~\GeV$ and $\mT>600~\GeV$ retain the hard tail while reducing the destructive contribution.

At $\MU=1.5~\TeV$, the magnitude of the interference contribution in $\bTagSR$-$\Res$ is typically below $5\%$ of the pure BSM yield over the parameter points studied.
In $\bTagSR$-$\NonRes$, fractions as large as $40\%$ are found at the lowest $\gU$ values, where the signal yield is too small to be excluded by the analysis.
Near the expected exclusion boundary from $\bTagSR$-$\NonRes$ alone, the contribution is typically below $10\%$ at $\MU=1.5~\TeV$ and below $5\%$ at $\MU=3.0~\TeV$.
The expected limits obtained with and without interference are found to differ little in the coupling range probed by $\bTagSR$.
These results support using the pure BSM yield to rank the trial selections in Section~\ref{subsec:ch7_scan}.
Interference is included in the nominal left-handed interpretation described in Section~\ref{sec:ch10_statistical_model}.

In $\bVetoSR$, the interference contribution is larger and is evaluated at each signal parameter point, as discussed in Section~\ref{sec:ch7_sr0b}.
The supplementary $\met$ scan in Appendix~\ref{app:sr0b_optimisation} illustrates how a tighter requirement can reduce this contribution in $\bVetoSR$.
